% !PRINTED 214
\weberpage{214}
%<<CONT
Equation (9), which is an identity with respect to $\tau$, remains
valid, however, when after carrying out the indicated raising to
powers $\tau^k$ is replaced on the right by $\tau_k$, and it therefore teaches
that between the functions $Y(t, x)$, $H(\tau, \xi)$ the relation holds

\weq{12}{H(\tau, \xi) = r\, Y(t, x),}

if we let the $t$ depend on the $\tau$ in the following way:
Set

\weq{13}{\begin{aligned}
t_{n-2} &= (\alpha\tau+\beta)^{n-2}\\
t_{n-3} &= (\alpha\tau+\beta)^{n-3}(\gamma\tau+\delta)\\
&\ \cdots\cdots\cdots\cdots\cdots\cdots\cdots\cdots\cdots\\
t_0 &= (\gamma\tau+\delta)^{n-2},
\end{aligned}}

and \gesperrt{after carrying out the powers} replace $\tau^k$ by $\tau_k$.

Now let $z$ be an arbitrary variable, and before we carry out
the raising to powers in (13), let us multiply these equations
in order by

\[
1, \quad -(n-2)z, \quad \frac{(n-2)(n-3)}{1\,.\,2}z^2, \cdots \pm z^{n-2};
\]

then on the left we obtain an integral rational function of $z$
of degree $(n-2)$

\weq{14}{\begin{aligned}
T(z) &= t_{n-2} - (n-2)t_{n-3}z\\
&\quad + \frac{(n-2)(n-3)}{1\,.\,2}t_{n-4}z^2 - \cdots \pm t_0 z^{n-2},
\end{aligned}}

and the right side yields, by the binomial theorem,

\[
[(\alpha\tau+\beta) - z(\gamma\tau+\delta)]^{n-2} = [\tau(\alpha-\gamma z) + (\beta-\delta z)]^{n-2},
\]

which, if we set

\weq{15}{z = \frac{\alpha\xi+\beta}{\gamma\xi+\delta}, \qquad \xi = \frac{\delta z-\beta}{-\gamma z+\alpha}}

goes over into

\[
(\alpha-\gamma z)^{n-2}(\tau-\xi)^{n-2} = \frac{r^{n-2}(\tau-\xi)^{n-2}}{(\gamma\xi+\delta)^{n-2}}
\]

If we carry out the power $(\tau-\xi)^{n-2}$ here by the binomial
theorem, and then replace $\tau^k$ by $\tau_k$, we obtain a transformation
of the function $T(z)$. We therefore set

% [sic] The print has a plus sign before \cdots; the binomial expansion and (14) require a minus sign there.

\[
\begin{aligned}
\Theta(\xi) &= \tau_{n-2} - (n-2)\tau_{n-3}\xi\\
&\quad + \frac{(n-2)(n-3)}{1\,.\,2}\tau_{n-4}\xi^2 + \cdots \pm \tau_0\xi^{n-2}
\end{aligned}
\]

and obtain the identical transformation
%CONT>>
